\section{问题建模}

\subsection{巡测场景与卫星遥感数据}

东海赤潮应急巡测场景建立在八年真实卫星遥感数据之上，真实海洋由任务当天卫星图构建，而航行器出发时仅携带3天前旧图。数据集采用美国国家海洋和大气管理局海岸观测网2018年1月1日至2026年10月23日共3197天的$2\text{ km}$逐日多星融合叶绿素a产品（经经验正交函数插补去云\cite{beckers2003dineof}，覆盖$25^\circ\text{--}33^\circ\text{N}$、$119.5^\circ\text{--}127^\circ\text{E}$，每日$385\times 361$格）及喷气推进实验室多尺度融合海表温度产品（重采样至$0.04^\circ$）\cite{chin2017mur}。全文除特别说明外，叶绿素a浓度、误差与遗憾均以$\log_{10}(\text{浓度}/\text{mg m}^{-3})$为单位。

我们在每年4至6月浙江外海陆架检索离岸至少$30\text{ km}$、平滑峰值超$5\text{ mg m}^{-3}$的藻华中心，随机平移至多$16\text{ km}$后截取$80\text{ km}\times 80\text{ km}$窗口（距岸$15\text{ km}$内禁航，可航行占比不低于80\%，窗内极差不低于0.3），共得2018至2022年开发集窗口274个、2023至2026年测试集窗口148个，航行器从窗口边缘随机可航行格入水（表\ref{tab:settings}）。鉴于东海表层赤潮斑块主要随水体与潮汐整体往复漂移\cite{lou2014diurnal}，本文在随表层水体漂移的参考系中建模，斑块相对水体静止，平流与生消作用体现在3天旧图与当天卫星图的实测演变统计中。

实验设置两种真值模型：主设置“当天卫星图 + 小尺度变化”在当天卫星图上叠加标准差$s_{\text{sat}}=0.15$、相关长度3格（$6\text{ km}$）的零均值高斯随机场，表征遥感平滑的亚像素斑块与反演残差（幅值按海洋一号丙星水色传感器在东海与南海实测验证约39\%的无偏相对差异即对数空间约0.15标定\cite{ye2021cocts}）；“当天卫星图”设置直接以当天卫星图本身为真值（$s_{\text{sat}}=0$）。两种设置下先验底图均为3天前卫星图$c_{\text{map}}(\mathbf{x})$。

\subsection{航行器动作空间与能耗模型}

参照百米级轻型航行器参数\cite{remus100}，航行器巡航速度$1.5\text{ m/s}$，电池容量$1\text{ kWh}$，巡航总功率$45\text{ W}$（含导航、传感器与低功耗处理器的平台功率$15\text{ W}$，推进功率$30\text{ W}$）。高功耗芯片参照嵌入式边缘图形处理器设定\cite{jetsonorinnx}，开启深想额外消耗$+20\text{ W}$（“省电加速器”为$+5\text{ W}$），推演速度为低功耗处理器的50倍。因二维双航路点规划器是对三维流场与避障路径评估的简化替代，实验以低功耗处理器每分钟推演次数表征基础算力（5、20或100次/分，扫描5至200次/分）。

在$T=24\text{ h}$内按$\Delta t = 20\text{ min}$划分为72个周期，机体动作与大脑规划并行执行\cite{xiao2020thinking, elboher2023formal}。机体选择“走”时沿当前航线航行$1.8\text{ km}$、每千米采样一次（24小时航程约$120\text{--}130\text{ km}$）；选择“停下”时关停推进器随水漂流，每4分钟对同一水体采样一次（每周期5次）。荧光计观测$y$与所在网格真实浓度$c(\mathbf{x})$满足：
\begin{equation}
  y = c(\mathbf{x}) + \varepsilon, \quad \varepsilon \sim \mathcal{N}(0, 0.1^2)
  \label{eq:obs}
\end{equation}
大脑规划可维持低功耗处理器后台规划，或开启高功耗芯片“深想”（推演次数扩至50倍）。新航线以当前周期结束时刻的预计位置为起点搜索，并在周期末替换旧航线；满24小时或电量耗尽即停止机动与采样。

\subsection{实时报告与时间平均遗憾}

在$T=24\text{ h}$内任意时刻$t \in [0, T]$，航行器从已访网格集合$\mathcal{V}(t)$中报告当前内部集合均值$\bar{c}_t(\mathbf{x})$最高的位置$\hat{\mathbf{x}}(t)$（若提前耗尽电量则冻结至$T$）：
\begin{equation}
  \hat{\mathbf{x}}(t) = \arg\max_{\mathbf{x}\in\mathcal{V}(t)} \bar{c}_t(\mathbf{x})
  \label{eq:answer}
\end{equation}
单次任务的时间平均遗憾$R$定义为全窗口真实最高浓度与实时报告位置真实浓度之差的时间均值：
\begin{equation}
  R = \frac{1}{T}\int_0^T \left[\max_{\mathbf{x}} c(\mathbf{x}) - c\big(\hat{\mathbf{x}}(t)\big)\right] dt
  \label{eq:regret}
\end{equation}
$R$越低表明越早锁定高浓核心区并保持足够续航；对数尺度下$R=0.1$和$0.3$分别对应报告浓度比真实峰值低约20\%和约一半。

\subsection{问题剖析}

内部计算与外部采样共享有限电池与时间窗口，存在三重矛盾。第一，增加计算的收益只能在内部想象海洋中估算，但规划器在搜索样本上的自评得分存在正向偏差：当推演次数扩大50倍时，优化器会微调航路点以穿过搜索样本中的偶然高值格，使内部得分虚高而实际航线不再改善\cite{smith2006curse, vanhasselt2010double}，因此需要将搜索航线与评判价值的想象海洋分离。

第二，观测的信息收益取决于3天前旧图的误差分布，遥感大数据虽能标定这一统计规律，却难以靠数值外推直接消除小尺度误差。正如第5.2节所示，深度学习时空预测无法从旧图中恢复数日后的小尺度峰值，航行器必须现场采样，并利用历史数据标定的误差方差与相关尺度构建先验想象海洋，随采样同化收缩不确定性。

第三，最优用时策略随底图不确定性、处理器速度与芯片功耗而变化，但任务间天然环境波动远大于策略间平均差异。在慢计算机上仅靠后台规划易使航行器在窗口中部局部绕圈，而在快计算机上开启高功耗芯片的额外能耗换取的航线改善十分有限。因此，无论是在下水前利用留出海洋开展全任务预演，还是在航行中利用留出海洋评估单步收益，核心均在于不依赖历史巡测任务而稳健剔除劣质策略。
